% ============================================================
%  APPENDIX — EXPERIMENT 2 FULL DETAILS  (news-response design, 2026-06-18)
%  Place after the main body, alongside the Experiment 1 appendix.
% ============================================================

\section{Experiment 2: Full Experimental Details}
\label{app:legacy-hidden-structure-details}

\subsection{Generative Process}
\label{app:exp2-dgp}
\begin{equation}
 o_t = o_0 + \phi(o_{t-1}-o_0) + g x_t + \epsilon_t, \qquad p_t=o_t+\eta_t.
\label{eq:exp2-dgp}
\end{equation}

Each episode is one city with a hidden news-responsiveness $g$, drawn once and never
revealed.  The latent opinion mean-reverts toward $o_0=50$ at rate $\phi$ while responding to
news in proportion to $g$, and is measured by a weekly poll as in Eq.~\eqref{eq:exp2-dgp};
opinion is clipped to $[0,100]$ and the poll is rounded to an integer.  The mean-reversion keeps
opinion inside the observable band, so the gain stays identifiable at every horizon rather than
saturating at a boundary.  Parameters:

\begin{table}[h]
\centering
\caption{Generative parameters for the news-response world.}
\label{tab:exp2-params}
\small
\begin{tabular}{ll}
\toprule
\textbf{Quantity} & \textbf{Value} \\
\midrule
News-responsiveness $g$        & $\sim\mathrm{Uniform}(0.1,\,1.0)$, one per city \\
Initial opinion $o_0$          & $50$ \\
Mean-reversion $\phi$          & $0.90$ \\
Horizon $K$                    & $10$ weeks (recovery read at each prefix $k=0,\dots,10$) \\
Weekly news dose $x_t$         & signed integer, $|x_t|\sim\mathrm{Uniform}\{3,\dots,12\}$, sign $\pm$ w.p.\ $\tfrac12$ \\
Test shocks                    & $\{-10,-5,+5,+10\}$ (four-point recovery slope) \\
Opinion noise $\sigma_o$       & $1.0$ \\
Poll noise $\sigma_p$          & $2.0$ \\
Cities per model              & $250$ (shared seeds across all forecasters) \\
\bottomrule
\end{tabular}
\end{table}

The news dose is shown to the agent as an explicit signed number, so the
news$\rightarrow$poll relationship is a controlled, visible quantity; the only thing hidden
is the gain $g$ that scales it.  Recovery is read at every prefix length $k=0,\dots,10$ by probing
the four test shocks $\{-10,-5,+5,+10\}$ and taking the regression slope of
Eq.~\eqref{eq:exp2-ghat}; the large shock magnitudes keep the estimate well-conditioned and the
four-point fit removes level-anchoring noise (a single shock conflates the two; see
Appendix~\ref{app:exp2-metrics}).

\begin{equation}
\hat g_k = \frac{\operatorname{Cov}_s(x^{(s)},\hat p_k^{(s)}-p_k)}{\operatorname{Var}_s(x^{(s)})}.
\label{eq:exp2-ghat}
\end{equation}

\subsection{The Bayesian Oracle}
\label{app:exp2-oracle}

The optimal forecaster marginalises the latent $g$ over a uniform grid of $41$ values on
$[0.1,1.0]$.  Under each hypothesis $g$, opinion is a linear-Gaussian state with known input
$g x_t$, so a Kalman filter on the observed polls gives both the exact log-likelihood and the
filtered opinion.  With state mean/variance $(\mu,v)$ initialised at $(50, 5^2)$, each step
updates
\begin{equation}
\mu^- = o_0 + \phi(\mu - o_0) + g x_t,\quad v^- = \phi^2 v + \sigma_o^2,\quad
s = v^- + \sigma_p^2,\quad
\ell \mathrel{+}= -\tfrac12\!\left(\log 2\pi s + \tfrac{(p_t-\mu^-)^2}{s}\right),
\end{equation}
followed by the Kalman correction $K=v^-/s$, $\mu=\mu^-+K(p_t-\mu^-)$, $v=(1-K)v^-$.  The
posterior over the grid is $P(g\mid p_{1:T})\propto \exp(\ell_g)$ (uniform prior); the
recovered $\hat g_{\mathrm{Bayes}}$ is its mean, and the predicted poll is the posterior-weighted
$\sum_g P(g\mid p_{1:T})\,\big(o_0+\phi(\mu_T^{(g)}-o_0) + g\,x_{T+1}\big)$.  Empirically this
estimator's recovery rises monotonically with the number of observed weeks, from
$\rho\approx0.67$ at $k=1$ to $\rho\approx0.94$ at $k=10$, with the slight pull-to-centre at the
extremes expected of a Bayesian estimator under finite data.

The two fixed shortcuts assume a constant gain and never update it: \emph{face value} uses
$g{=}1$ (predict $p_T + x_{T+1}$) and \emph{ignore news} uses $g{=}0$ (predict $p_T$).

\subsection{Secondary Informativeness References: Same-Information Forecasters}
\label{app:exp2-parity-refs}

The primary fair bar is the naive frequentist on the \emph{error} axis (\S\ref{app:exp2-metrics}).
As a secondary reference on the \emph{informativeness} ($\rho$) axis we additionally construct two
same-information forecasters handed \emph{exactly the same brief and nothing more} as the model,
both read through the identical four-shock predict-then-slope pipeline.  These characterize how much
of the across-city ordering is recoverable by a principled estimator told the prompt's structure;
one caveat, below, is that under rank correlation a one-observation estimate is unidentified and so
scores $\rho{=}0$ at $k{=}1$---a property of the metric at $k{=}1$, not a claim that the agents beat
a well-posed estimator (which, on error, they do not).  Each forecaster is restricted to information
the prompt actually gives:

\begin{itemize}
\item \textbf{No known baseline, no mean-reversion.}  The prompt never states $o_0{=}50$ and never
mentions $\phi$, so the forecasters use neither; they regress only the \emph{consecutive} weekly
poll changes $\Delta p_t = p_t - p_{t-1}$ on the news dose $x_t$.  The omitted mean-reversion pull
$-(1-\phi)(o_{t-1}-o_0)$ in $\Delta p_t$ is a shared handicap the model also pays.
\item \textbf{Believe the prompt (zero intercept).}  The prompt states the relationship is
$\Delta o\approx g\,x$ plus a \emph{zero-mean} drift, i.e.\ no systematic intercept, so the
forecasters take it at its word and fit through the origin.  A single change then identifies $g$, so
the first estimate is at $k{=}2$.  This is the parity analogue of what the model is told; a
forecaster that instead estimated a free intercept (the baseline reference below) would discard that
disclosed structure.
\item \textbf{A value at every $k$.}  The model is probed at every prefix, so the references answer
at every prefix too: at $k=0,1$, where no change is yet available, each returns the range midpoint
($0.55$)---a constant, hence recovery $\rho=0$.  These are real plotted points (``no recovery yet''),
not gaps.  Note even the midpoint cannot help at $k{=}1$: with one poll and no baseline $g$ is
unidentifiable, so any baseline-free forecaster scores $\rho\approx0$ there while the oracle, handed
$o_0{=}50$, already reaches $0.66$---the gap is exactly the value of the known baseline.
\item \textbf{Stated range and noise only.}  The $g\in[0.1,1.0]$ range and the noise scales are
stated, so the forecasters may use them: the OLS slope is clamped to $[0.1,1.0]$ (without the clamp
$\approx10\%$ of horizon estimates fall outside it), and the Bayes likelihood uses
$\sigma=\sqrt{\sigma_o^2+2\sigma_p^2}=3$, the value implied by the prompt's ``about 1''/``about
2.''  Recovery is essentially invariant to this scale (sweeping $\sigma\in[1,8]$ moves $\rho$ by
$\leq0.02$), so it confers no hidden edge.
\end{itemize}

The \emph{OLS forecaster} is the through-origin slope $\hat g=\sum_t x_t\,\Delta p_t/\sum_t x_t^2$
(clamped); the \emph{Bayes forecaster} is the posterior mean of $g$ over the uniform grid under
$\Delta p_t\sim\mathcal N(g\,x_t,\sigma^2)$.  Both are well calibrated at the horizon (mean
$\overline{\hat g}\approx0.55$, matching the true mean; regression slope $\beta\approx0.90$ for OLS
and $0.80$ for Bayes, the lower value reflecting benign prior shrinkage), and the forecast wrapper is
a no-op here---clipping rarely binds, so the recovered slope equals the direct estimate.  The two are
nearly identical---$\rho\approx0.51$/$0.54$ at $k{=}2$ rising to $\approx0.89$ at $k{=}10$---and sit
below the oracle most where data is scarce ($0.72$ at $k{=}2$, the value of the privileged baseline
and $\phi$) and close by the horizon ($0.94$).  Each slope panel also carries the matching
\emph{oracle} as a faint dotted peak: the Bayes oracle in the Bayes panel and, in the OLS panel, an
\emph{OLS oracle}---the frequentist analogue that likewise knows $o_0{=}50$ and $\phi$ and regresses
the $\phi$-corrected change $y_t=(p_t-o_0)-\phi(p_{t-1}-o_0)$ on $x_t$ through the origin.  Knowing
the baseline, it identifies $g$ from $k{=}1$ ($\rho\approx0.62$) and rises to $0.92$, tracking just
under the Bayes oracle (no prior or optimal filtering).  The oracles are plotted only as faint
ceilings; the fair verdict is a model's position relative to the two forecasters, which we mark by
shading each slope panel green above its forecaster (the model infers more than a principled
same-information fit) and red below.

\paragraph{Relation to the baseline naive-OLS.}  Both conditions use the same data-only forecasters
with the same convention---consecutive poll changes, no $o_0{=}50$ anchor (the baseline reference,
Appendix~\ref{app:exp2-metrics}, also drops it)---so the references are a consistent yardstick and
the parity manipulation changes only what the \emph{model} is told.  The one principled difference is
the intercept, and it tracks the disclosure: the \emph{parity} forecaster believes the stated
zero-intercept structure and fits through the origin, identifying $g$ from a single change
($\rho\approx0.51$ at $k{=}2$), whereas the \emph{baseline} forecaster---whose prompt never states
the structure---must fit a free intercept and is under-determined with few weeks
($\rho\approx0.00$ at $k{=}2$, $0.13$ at $k{=}3$).  The two converge by the horizon
($\approx0.86$--$0.89$); the early gap is precisely the sample-efficiency value of disclosing the
generative structure---the forecaster analogue of the experiment's central effect.

\subsection{Prompt (multi-shock reasoning probe)}
\label{app:exp2-prompt}

At each prefix length the agent is shown the observed weeks as a table, asked to first
\emph{estimate} the city's per-news response and then predict, and queried once at each of the four
test shocks $\{-10,-5,+5,+10\}$.  Its recovered gain $\hat g_k$ is the slope of
Eq.~\eqref{eq:exp2-ghat} across those four predictions, and we also store the stated rate.
Responses are JSON; we accept only a well-formed object containing \texttt{predicted\_poll} and
treat a missing or truncated object as missing data (with no fallback to stray numbers from the
reasoning text), retrying up to three times before recording a miss.

\begin{small}
\begin{verbatim}
You are tracking a local election campaign with a weekly tracking poll.
Each week reports the net news (a signed number: positive = good news for
the candidate, negative = bad) and the resulting poll (support, 0-100).

How this city works (the same for every week):
  - This city has a fixed news-responsiveness g, a number between 0.1 and
    1.0 (you do not know which).
  - Each week the underlying opinion moves by about g times that week's net
    news, plus a small random drift of about 1 point.
  - The reported poll is that opinion plus survey noise of about 2 points.
  - Low g = sticky coverage (news barely moves the poll); high g = news
    moves it nearly one-for-one.

| Week | Net news | Poll |
|------|----------|------|
|    1 |    -3 |   49 |
|    2 |    -4 |   45 |
|    3 |    +9 |   54 |
|    4 |   +12 |   64 |

From the history above, estimate this city's g (poll points moved per +1
point of net news).

Given that next week's net news will be <shock>, predict next week's poll.
You may reason step by step. Your reply is capped at about 500 tokens,
which is ample room: use as much or as little reasoning as the inference
needs, and be sure to finish the JSON within the limit.
END your reply with the answer as a JSON object on its own line:
{"rationale": "one short sentence", "points_per_news": <number>,
 "predicted_poll": <integer 0-100>}
\end{verbatim}
\end{small}

The open-weight models (Qwen~2.5-7B/14B/32B/72B, Llama~3.1-8B) are served via Ollama at temperature
$0.1$; the frontier agents (GPT-5.4, Claude~Opus~4.8, DeepSeek~V4-Pro) are queried through their
APIs at the same temperature.  All models reason step by step before answering, under a $2048$-token
completion budget, with the prompt itself instructing brevity (``capped at about 500 tokens \ldots a
few short sentences'')---enough for the most verbose models to finish (at smaller budgets the largest
models truncate mid-reasoning, which the hardened parser records as a miss rather than a spurious
number).  The exact per-run configuration---\texttt{max\_tokens}, \texttt{reason\_cap}, and a hash of
the prompt template (\texttt{prompt\_sha})---is recorded in each run's manifest, so the prompt and
budget for any result file are always recoverable.

\subsection{Recovery and Cost Metrics}
\label{app:exp2-metrics}

Over cities $\{i\}$ with true gains $g_i$ and recovered gains $\hat g_{i,k}$
(Eq.~\eqref{eq:exp2-ghat}, after $k$ observed weeks) we lead with two per-instance \emph{error}
metrics.  The first is the \emph{recovery error} $\varepsilon_k=\frac1N\sum_i|\hat g_{i,k}-g_i|$
(lower is better): did the forecaster infer the hidden gain?  The second is the \emph{next-poll
forecast error} $\pi_k=\frac1N\sum_i\frac14\sum_{s}\bigl|\hat p_{i,k}(s)-\mu_{i,k}(s)\bigr|$, the
mean gap (in poll points) between the poll a forecaster predicts under each held-out shock
$s\in\{-10,-5,+5,+10\}$ and the ground-truth expected poll
$\mu_{i,k}(s)=\mathrm{clip}\bigl(50+\phi(o_{i,k}-50)+g_i s,\,0,\,100\bigr)$ that shock produces in the
generative process (with $o_{i,k}$ the true latent opinion after $k$ weeks).  Because the four shocks
are counterfactual they never realize, so the target is the conditional \emph{mean}---the
Bayes-optimal point forecast, which excludes the irreducible weekly drift and survey noise---and even
the oracle, which must filter $o_{i,k}$ and $g_i$ from the polls, sits above zero.  We pair these with
the \emph{informativeness} $\rho_k=\rho(\hat g_k,g)$ as a distinct, secondary axis, and additionally
report the regression slope $\beta=\mathrm{Cov}(g,\hat g)/\mathrm{Var}(g)$ ($\beta\!\to\!1$ tracks the
latent, $\beta\!\to\!0$ ignores it); the mean $\overline{\hat g}$ (a magnitude/calibration check); and
the correlation of the model's \emph{stated} rate with $g$.  The recovery error is summarised in the
main-text Table~\ref{tab:exp2-skill}; the next-poll error is tabulated per model in
Table~\ref{tab:exp2-poll}.

\paragraph{Why error, not correlation, is the headline.}  The claim under test---a useful prior yields
a sensible estimate from sparse evidence (the carnival coin)---is about the \emph{calibration} of a
single estimate, not about ordering a population.  Rank correlation is the wrong headline for it
because $\rho$ is scale- and location-invariant: a forecaster whose $\hat g$ is an affine image of
$g$, $\hat g=a+bg$ with any $b>0$, attains $\rho=1$ regardless of $a$ and $b$.  An unregularized
regression on one or two noisy weeks does exactly this---it preserves the across-city ordering while
badly \emph{overshooting} the magnitude ($b\gg1$), so it scores high $\rho$ while being poorly
calibrated ($\varepsilon$ larger than simply guessing the prior mean; see the naive frequentist at
$k{=}2$ in Table~\ref{tab:exp2-skill}).  A prior-regularized forecaster makes the opposite trade,
shrinking toward the range and lowering $\varepsilon$ at a small cost in early $\rho$.  Reporting only
$\rho$ would therefore \emph{reward} the overshoot that a good world model avoids; we lead with
$\varepsilon$ and read $\rho$ alongside it as the informativeness axis, and a forecaster is credited
with recovering structure only when it is strong on \emph{both}.

\paragraph{The fair statistical bar: a naive frequentist.}  The primary reference for
\emph{calibration} is a \emph{naive frequentist}: an unregularized through-origin regression of the
consecutive poll changes $p_t-p_{t-1}$ on the news dose $x_t$ over the first $k$ weeks, using only
the data the agent sees---no prior, no clamp, no noise scales or $\phi$.  It needs two observations
to be identified, so at $k{=}1$ the honest floor is instead guessing the prior mean ($0.55$).  This
reference makes the sparse-regime point sharply on the error axis: with one or two noisy weeks it
\emph{overshoots}, so its error is actually \emph{worse} than guessing the prior at $k{=}2$
($\varepsilon\approx0.35$ vs.\ $0.22$) before settling toward the oracle as data accumulates
($\varepsilon\approx0.16$ by $k{=}5$).  An agent that beats it in the sparse regime is regularizing
where a one-line fit cannot.  We also report, as a secondary \emph{informativeness} reference, the
same-information Bayes and OLS \emph{forecasters} of Appendix~\ref{app:exp2-parity-refs}, which take
the prompt's disclosed structure at its word (through-origin fit); these rank-order cities well but,
being data-only, share the frequentist's early miscalibration.

\paragraph{Why multiple shocks.}  A single-shock readout, $\hat g=(\hat p-p_T)/x$, is
contaminated: writing the prediction as $\hat p = \hat o_T + \hat g\,x$, the implied gain becomes
$\hat g + (\hat o_T - p_T)/x$, so any error in where the forecaster anchors the level enters
divided by $x$.  Regressing across the four shocks (Eq.~\eqref{eq:exp2-ghat}) sends that level term
into the intercept and leaves the slope as a clean estimate of the gain.  The effect is
substantial: a single-shock read materially lowers recovery for every forecaster, including the
oracle, and can invert the model ordering by conflating level anchoring with the gain.  We therefore
treat the four-shock slope as the faithful measure.

\subsection{Extended Results}
\label{app:exp2-extended}

Table~\ref{tab:exp2-skill} and Figure~\ref{fig:exp2-skill} (main text) give the skill/calibration
summary.  Three points warrant emphasis, now on the informativeness axis $\rho$ (which the main text
demotes to secondary; recall $\rho$ alone rewards the overshoot a good prior avoids,
\S\ref{app:exp2-metrics}).  First, $\rho$-recovery rises with the number of
observed weeks for every model---they accumulate evidence rather than applying a fixed gain---and the
asymptotic level orders by capability: frontier GPT-5.4/Claude $\approx0.84$/$0.82$, open-weight
Qwen~2.5-32B/72B $\approx0.59$, down to Qwen~2.5-7B $0.29$ and Llama~3.1-8B $0.19$ (parity, $k=10$).
Second, on the \emph{error} axis the fair bar is the naive frequentist, not the oracle: the
frontier agents beat it in the sparse regime---where its unregularized fit overshoots---and converge
with it as data accumulates, so an agent should be judged by whether it stays calibrated where the
statistical fit cannot, and the strongest agents do (GPT/Claude
beat it through $k\approx4$) but plateau and are overtaken by the horizon, while open-weight models
trail it throughout.  Third,
recovery does not entail calibration: several models track $g$ in rank while biasing $\overline{\hat
g}$ high (Llama~3.1-8B $0.79$, GPT-5.4 $0.66$ vs.\ truth $0.55$), and rank and slope can diverge
(Qwen~2.5-7B has higher $\beta$ but lower $\rho$ than 72B), so structural recovery and forecast
calibration are distinct.

\paragraph{Next-poll forecast error (the second headline metric).}  Table~\ref{tab:exp2-poll} reads
the same probe through the forecast the recovered gain feeds: the mean gap between the poll each
forecaster predicts under the four held-out shocks and the ground-truth expected poll $\mu_k(s)$
(\S\ref{app:exp2-metrics}).  The pattern mirrors the recovery axis.  The frontier agents forecast to
$\approx2.0$--$2.4$ poll points, approaching the oracle ($1.2$--$1.9$) and far below every fixed
shortcut (face value $\approx3.9$, ignore-news $\approx4.3$); the naive frequentist \emph{overshoots}
the same way it does on $\hat g$, forecasting worse than guessing the prior at $k{=}2$ ($3.1$ vs.\
$2.3$), and the agents beat it through the sparse regime before pulling ahead of the always-prior
forecast by $k{=}5$ ($2.0$ vs.\ $2.4$).  The green/red shading is against the frequentist bar (the
always-prior floor at $k{=}1$), which converges quickly once identified, so the smaller open-weight
models---strong against the shortcuts but weaker on level---turn red at larger $k$; the ordering by
scale is monotone (Llama~3.1-8B, near face-value, forecasts at $6$--$10$ points).  Because the metric
now includes the forecast \emph{level}, not just the slope, it is a stricter test than $\hat g$
recovery, and the agents remain the only non-oracle forecaster near the ceiling throughout.

% ----------------------------------------------------------------------------
% Static copy of the generated table. SOURCE OF TRUTH: paper/generate_exp2_poll_table.py
% To refresh: run `python paper/generate_exp2_poll_table.py` and paste tables/exp2_poll_table.tex.
% Requires booktabs and \usepackage[table]{xcolor}.
% ----------------------------------------------------------------------------
\begin{table}[t]
\centering
\caption{\textbf{Next-poll forecast error} $\overline{|\hat p_{t+1}-p_{t+1}|}$ (poll points, lower is better), the second headline metric: does recovering $g$ translate into good forecasts? For each held-out shock $s$ the target is the ground-truth expected poll $\mu_k(s){=}\mathrm{clip}(50+\phi(o_k{-}50)+g s)$; the counterfactual shocks never realize, so this is the DGP's conditional mean, the Bayes-optimal point forecast (even the oracle, which must filter $o_k$ and $g$, sits above $0$). The \emph{Bayes oracle} is the ceiling and the fixed shortcuts (face value $g{=}1$, ignore news $g{=}0$) the floor. The fair statistical bar is the \emph{naive frequentist} (unregularized slope, predict $p_T{+}\hat g s$); undefined at $k{=}1$, so there the bar is \emph{always-guess-the-prior}. Cells are \colorbox{green!22}{green} where a model beats that per-$k$ bar, \colorbox{red!18}{red} where it falls short. The frontier agents forecast to $\approx2$ poll points---near the oracle, far below every shortcut, and beating the overshooting frequentist through the sparse regime. Frontier $n{=}100$; open-weight $n{=}68\text{--}250$.  \textsuperscript{$\dagger$}Frontier probe.}
\label{tab:exp2-poll}
\footnotesize
\begin{tabular}{l@{\hskip 6pt}cccc}
\toprule
\textbf{Forecaster} & $k{=}1$ & $k{=}2$ & $k{=}3$ & $k{=}5$ \\
\midrule
Bayes oracle & 1.91 & 1.60 & 1.50 & 1.19 \\
Naive freq. & --- & 3.13 & 2.55 & 2.15 \\
Always-prior & 2.39 & 2.31 & 2.46 & 2.39 \\
Face value ($g{=}1$) & 3.87 & 3.87 & 3.92 & 3.88 \\
Ignore news ($g{=}0$) & 4.27 & 4.20 & 4.27 & 4.24 \\
\midrule
Claude Opus~4.8\textsuperscript{$\dagger$} & \cellcolor{red!18}2.40 & \cellcolor{green!22}2.36 & \cellcolor{green!22}2.44 & \cellcolor{green!22}2.04 \\
GPT-5.4\textsuperscript{$\dagger$} & \cellcolor{green!22}2.36 & \cellcolor{green!22}2.45 & \cellcolor{green!22}2.41 & \cellcolor{green!22}2.04 \\
DeepSeek V4-Pro\textsuperscript{$\dagger$} & \cellcolor{red!18}2.54 & \cellcolor{green!22}2.67 & \cellcolor{green!22}2.35 & \cellcolor{green!22}2.07 \\
\midrule
Qwen~2.5-72B & \cellcolor{red!18}3.20 & \cellcolor{green!22}2.72 & \cellcolor{red!18}2.61 & \cellcolor{red!18}2.52 \\
Llama~3.3-70B & \cellcolor{red!18}2.72 & \cellcolor{green!22}3.11 & \cellcolor{red!18}2.76 & \cellcolor{red!18}2.53 \\
Llama~3.1-70B & \cellcolor{red!18}2.87 & \cellcolor{green!22}2.67 & \cellcolor{red!18}2.60 & \cellcolor{red!18}2.42 \\
Qwen~2.5-32B & \cellcolor{red!18}2.99 & \cellcolor{green!22}2.54 & \cellcolor{green!22}2.48 & \cellcolor{red!18}2.26 \\
Mistral-Small-24B & \cellcolor{red!18}3.08 & \cellcolor{green!22}2.67 & \cellcolor{red!18}2.63 & \cellcolor{red!18}2.37 \\
Qwen~2.5-14B & \cellcolor{red!18}3.33 & \cellcolor{green!22}2.85 & \cellcolor{red!18}2.68 & \cellcolor{red!18}2.53 \\
Mistral-Nemo-12B & \cellcolor{red!18}4.52 & \cellcolor{red!18}3.84 & \cellcolor{red!18}4.35 & \cellcolor{red!18}5.55 \\
Llama~3.1-8B & \cellcolor{red!18}10.31 & \cellcolor{red!18}7.14 & \cellcolor{red!18}6.30 & \cellcolor{red!18}6.86 \\
Qwen~2.5-7B & \cellcolor{red!18}3.54 & \cellcolor{green!22}2.94 & \cellcolor{red!18}2.91 & \cellcolor{red!18}3.11 \\
\bottomrule
\end{tabular}
\end{table}

\begin{figure}[t]
\centering
\includegraphics[width=0.62\linewidth]{figures/exp2_informativeness.pdf}
\caption{\textbf{Calibration vs.\ informativeness at the sparse anchor $k{=}2$} (the view the main-text
figure demotes in favour of the two error axes).  $x$ = informativeness $\rho(\hat g,g)$ (higher is
better ordering), $y$ = calibration error $\overline{|\hat g-g|}$ (lower is better).  The fixed
shortcuts sit in the poor corners; the naive frequentist buys $\rho$ at the cost of calibration
(upper-right); always-prior is calibrated but blind ($\rho{=}0$).  Only the frontier agents reach the
good corner (bottom-right), informative \emph{and} calibrated---the joint property a single ordering
number would hide.}
\label{fig:exp2-informativeness}
\end{figure}

\subsection{Data-Quality and Robustness Checks}
\label{app:exp2-robustness}

The headline---agents recover the latent quickly but \emph{plateau} (or even dip) while the oracle
and forecaster references keep accumulating---rests on the \emph{shape} of recovery at large $k$, so
we audit data quality per model and per prefix length, on both the baseline and information-parity
conditions.  The open-weight models run the full $250$ cities; the frontier agents (GPT-5.4,
Claude~Opus~4.8, DeepSeek~V4-Pro) run $n{=}100$, with stable per-$k$ estimates (bootstrap
SE${\approx}0.04$--$0.05$).

\textbf{Missing predictions.}  After up to three retries, the fraction of probes returning no
parseable forecast is $\approx0\%$ across models and $k$ (the hardened parser of
Appendix~\ref{app:exp2-prompt} records a truncated or malformed reply as missing rather than
scraping a number from the reasoning).  The sole exception is Qwen~2.5-7B in the parity condition at
$k=0$, where the ``no weeks observed yet'' prompt yields no JSON on $\approx46\%$ of probes; since
$k=0$ is the no-evidence prior (recovery $\approx0$ regardless) we report from $k\geq1$.  No
accumulation failure occurs at the horizons that carry the claim.

\textbf{Extreme predictions.}  Forecasts pinned to a boundary ($\hat p\leq2$ or $\geq98$) are
rare---$0.1$--$2.0\%$ of predictions overall, lowest under parity---and, crucially, do \emph{not}
increase with $k$: high-horizon ($k=8$--$10$) extreme rates are $\approx1$--$3\%$ (baseline) and
$\lesssim1\%$ (parity).  A small spike at $k=1$ (models over-reacting to a single week) does not
propagate.  High-$k$ recovery is therefore not depressed by boundary predictions.

\textbf{Response consistency.}  The one quality gradient is monotonicity: a faithful forecaster
should predict a weakly larger poll for a larger shock, yet some four-shock probes are
non-monotonic.  This is low and roughly flat in $k$ under parity (Qwen~2.5-72B $3$--$13\%$, 32B
$7$--$18\%$) but rises with $k$ under the baseline prompt, sharply for the small models (72B
$13\%\!\to\!44\%$ over $k$; 7B $54$--$66\%$ throughout).  Because a non-monotonic quadruple yields a
noisier slope (regression dilution), rising non-monotonicity could in principle attenuate recovery
at large $k$ and \emph{manufacture} a plateau.

\textbf{Ruling out the artifact.}  Two controls show the plateau is real recovery, not attenuation.
First, the plateau is equally pronounced in the \emph{parity} condition, where non-monotonicity is
low and flat: recovery rises to its level by $k\approx5$ and holds (Qwen~2.5-72B even dips, $0.65$ at
$k=5$ to $0.58$ at $k=10$) while the forecaster climbs past it ($0.79\!\to\!0.89$).  Second,
restricting the recovery correlation at each $k$ to the \emph{monotonic-only} cities leaves the
curves essentially unchanged in both conditions---the models still plateau and the forecaster still
overtakes them by the horizon.  We therefore read the non-monotonicity itself as a
\emph{secondary behavioral signal}---coherent shock-response degrades with context length,
especially for small and un-prompted models---rather than as a measurement confound.  We report the
parity condition as the primary (clean) measure and the baseline as a robustness-confirmed
secondary.

\subsection{Qualitative: over-identifying \texorpdfstring{$g$}{g} from a single week}
\label{app:exp2-traces}

The $k{=}1$ over-reaction noted above has a clean qualitative signature in the reasoning traces, and
inspecting it also explains an apparent ordering wrinkle: among the larger open-weight models
Qwen~2.5-72B trails 32B at $k{=}1$ ($\rho\,0.24$ vs.\ $0.45$ on the shared cities), though a paired
bootstrap puts the gap within noise at the current $n$ ($95\%$ CI on $\rho_{32}-\rho_{72}$ is
$[-0.02,+0.45]$, and the two are statistically tied by $k{=}5$). We queried the live Qwen~2.5-72B
server with the parity prompt at $k{=}1$ and logged the full generations; every reply completed
normally (\texttt{finish\_reason}{=}\texttt{stop}, no truncation, $\approx1$k characters), so the
behavior below is a genuine inferential choice rather than a cut-off, consistent with the
$\approx0\%$ missing-prediction rate of \S\ref{app:exp2-robustness}.

The model assumes the (undisclosed) baseline $o_0{=}50$ and then point-estimates $g$ by dividing a
single week's deviation from $50$ by that week's news---reading one noisy observation as clean
signal. For city~91 (true $g{=}0.18$, a very \emph{sticky} city), week~1 shows news $-5$ and poll
$47$:
\begin{quote}\small\itshape
``assuming initial opinion was around 50\,\ldots\ the change in opinion from 50 to 47 is approximately
$-3$. This gives us $-3 = g\times(-5)+1$\,\ldots\ $g = 0.8$.''
\end{quote}
The $-3$ is almost entirely drift and survey noise, but it is read as responsiveness, yielding
$\hat g{=}0.8$ for a city whose true $g$ is $0.18$. The model even flags the tension and overrides
its own caution---``$g = 1.0$. However, this seems high given the noise and drift\,\ldots\ a more
reasonable estimate\,\ldots\ is $g \approx 0.8$''---and still commits to the inflated value. City~7
($g{=}0.39$) is identical: ``$g = 5/5 = 1.0$\,\ldots\ let's use $0.9$.''

Because a single poll cannot separate $g$ from the baseline and the per-week noise (the
identifiability point of \S\ref{app:exp2-parity-refs}: with one observation $g$ is unidentified),
this confident point-estimate over-fits the noise---worst on low-$g$ cities, where the true signal is
smallest relative to it. It is a coherent but over-eager inference, not a hedge toward the prior; the
smaller Qwen~2.5-32B is empirically more regularized at $k{=}1$, which is why its early $\hat g$
tracks the true $g$ better. As $k$ grows the deviations accumulate, the noise averages down, and the
two converge. The episode is a small illustration of the section's theme: in the sparse regime the
useful move is to \emph{withhold} a confident estimate, and more scale does not automatically buy
that restraint.

\subsection{Prompt sensitivity of the few-shot edge}
\label{app:exp2-promptsens}

The over-identification of \S\ref{app:exp2-traces} raises a methodological question: is the early
($k{=}1$) recovery a property of the agent, or an artifact of how we elicit $\hat g$? We probed this
by varying \emph{only} the clause of the prompt (\S\ref{app:exp2-prompt}) that steers how the model
should treat a short, noisy history, holding the generative-structure disclosure fixed:
\emph{(i) assertive}---``estimate this city's $g$, seeing through the survey noise'' (implies the
noise is removable); \emph{(ii) hedging}---``if the weeks so far cannot determine $g$, keep your
estimate near the middle of its range'' (instructs regularization); and \emph{(iii) neutral}---
``estimate this city's $g$,'' with no instruction on handling the uncertainty.

The early readout moves with the steer while the rest of the curve does not. For the frontier agents
(the cleanest sample), $\rho_1$ is ${\approx}.30$ under the assertive prompt, falls toward $0$ under
the hedging prompt, and sits at ${\approx}.20$ under the neutral one, whereas $\rho_3$ and $\rho_5$
are essentially unchanged across all three (${\approx}.6$ and ${\approx}.75$). So the \emph{level} of
the sparse-data edge is prompt-sensitive, but its \emph{shape}---rising with $k$ toward the
forecaster---is not. We therefore adopt the \textbf{neutral} prompt as primary: the assertive wording
manufactures over-confidence (it asserts a signal the single week does not contain) and the hedging
wording manufactures restraint, so only the neutral readout reflects the model's \emph{own}
disposition. Under it, at $n{=}100$ per frontier model, GPT-5.4 and Claude retain a genuine but
\emph{modest} single-week edge ($\rho_1{\approx}.19$--$.22$, with calibration error
$\varepsilon_1{\approx}.20$ near the oracle's $.18$), climbing to $\rho_5{\approx}.76$ by $k{=}5$; the open-weight models do not recover from one week
($\rho_1{\approx}0$ or negative, e.g.\ Qwen~2.5-7B $-.10$, Llama~3.1-8B $-.26$). This
tempers---without overturning---the few-shot story: there is a real single-week advantage for the
strongest agents, but a smaller one than an assertive elicitation would suggest. (The three
conditions were sequential design iterations, not a matched same-$n$ sweep; we read the directional
sensitivity as motivation for the neutral choice and report the neutral condition at full $n$, with a
controlled three-arm comparison left to future work.)

\subsection{The carnival-coin vignette: do frontier agents actually give the world-model answer?}
\label{app:exp2-carnival}

The Introduction motivates the paper with a cartoon: an observer sees two heads in two flips of
an unfamiliar coin at a carnival and must forecast the next flip.  A \emph{na\"ive} forecaster
assumes a fair coin and answers $0.5$; a \emph{frequentist} uses only the two flips and answers
$1.0$; a \emph{world-model} forecaster reasons about an unknown, possibly-biased coin and lands in
between. Here we test the vignette directly with
the three frontier agents and record what they actually say.

\paragraph{Probe.}  We reuse the Experiment~2 evaluation harness (the same Azure AI Foundry clients
and call path as the recovery probe) and query each agent $n{=}8$ times at temperature $0.7$ on the
neutral prompt below.  Two design choices keep the elicitation fair.  First, the prompt establishes
only that the coin is \emph{unfamiliar} and that the two flips are the sole evidence; it never raises
the fair-versus-rigged framing, since naming either side pre-loads the answer (calling it simply ``a
coin'' cues the textbook fair-coin reflex, and an early pilot with that wording drew a flat $0.5$
from every model).  Second, we forbid calculation: the agent must reason in plain words rather than
recite a rule.  Without this constraint all three models collapse onto Laplace's rule of succession
and return an identical, rigid $3/4=0.75$; the no-math instruction elicits the graded
judgment the figure is about, not a formula lookup.

\begin{small}
\begin{verbatim}
You are at a carnival game. A coin you have never seen before is flipped
twice in front of you, and it comes up heads both times. These two flips
are everything you know about this coin.

What is the probability that the next flip comes up heads?

Think about it intuitively, the way you actually would in the moment -- do
NOT use any formula, rule, or calculation, and do not work out any numbers
in your reasoning. Just reason in plain words about what you'd expect and
why. Then END your reply with your gut answer as a JSON object on its own
line:
{"rationale": "one short sentence, no math",
 "p_heads": <number between 0 and 1>}
\end{verbatim}
\end{small}

\paragraph{Result.}  Every agent places the next flip above even odds but well short of certainty,
exactly the world-model region of the cartoon (Table~\ref{tab:carnival}).  Per-model medians span
$0.67$--$0.75$ and the pooled mean over all $24$ samples is $0.68$---closer to the figure's
world-model value ($\approx0.65$) than to either the na\"ive ($0.5$) or frequentist ($1.0$) anchor.
None of the $24$ completions reverted to $0.5$ or $1.0$.  Notably, Claude and DeepSeek frequently
cite the \emph{setting} unprompted---``a sketchy carnival,'' ``carnival coins are often rigged''---to
justify a head-leaning coin, importing world knowledge from the single word ``carnival'' with no
steer from us; GPT-5.4 stays more conservative, leaning on the thinness of two flips
(Table~\ref{tab:carnival-samples}).  This is the same disposition Experiment~2 measures with a poll
instead of a coin: a useful prior reads some signal from a single observation.

\begin{table}[h]
\centering
\caption{Frontier-agent forecasts on the carnival-coin vignette ($n{=}8$ samples each, temperature
$0.7$, no-math elicitation).  Reference anchors: na\"ive $0.5$, frequentist $1.0$, figure
world-model $\approx0.65$.}
\label{tab:carnival}
\small
\begin{tabular}{l c c c}
\toprule
Model & Median $P(\text{heads})$ & Mean & Range \\
\midrule
Claude Opus~4.8  & $0.675$ & $0.68$ & $[0.65,\,0.75]$ \\
GPT-5.4          & $0.670$ & $0.64$ & $[0.60,\,0.67]$ \\
DeepSeek V4-Pro  & $0.750$ & $0.71$ & $[0.60,\,0.75]$ \\
\midrule
Pooled ($n{=}24$) & $0.700$ & $0.68$ & $[0.60,\,0.75]$ \\
\bottomrule
\end{tabular}
\end{table}

\begin{table}[h]
\centering
\caption{Representative verbatim rationales (one short sentence each, as emitted).  The agents reason
qualitatively about an unknown, possibly-biased coin rather than applying a rule.}
\label{tab:carnival-samples}
\small
\begin{tabular}{l c p{8.3cm}}
\toprule
Model & $P$ & Rationale \\
\midrule
Claude Opus~4.8 & $0.75$ & Carnival coins are often rigged and two heads makes me suspect it favors heads. \\
Claude Opus~4.8 & $0.65$ & Two heads gently nudge me toward heads, and the carnival setting adds suspicion, but it's weak evidence. \\
GPT-5.4         & $0.67$ & Two heads in a row makes this unfamiliar coin feel somewhat more head-leaning, but not overwhelmingly so. \\
GPT-5.4         & $0.60$ & Two heads in a row makes heads feel somewhat more likely next, but the evidence is still very thin. \\
DeepSeek V4-Pro & $0.75$ & A carnival coin showing two heads seems likely rigged or biased toward heads. \\
DeepSeek V4-Pro & $0.60$ & Two heads in a row slightly favors the coin being biased toward heads, but the evidence is weak. \\
\bottomrule
\end{tabular}
\end{table}

The probe script (\texttt{eval/run\_coin\_probe.py}) and the raw per-sample log
(\texttt{data/coin\_probe/}) are included with the Experiment~2 code; the saved JSONL records the
exact prompt, sampling settings, and every completion so the table is reproducible.

\subsection{Limitations and Future Directions}
\label{app:exp2-limitations}

\textbf{Frontier-agent sample size.}  The three frontier agents (GPT-5.4, Claude~Opus~4.8,
DeepSeek~V4-Pro) are run on $n{=}100$ cities---a verified subset of the open-weight $250$ (seeds
$0\ldots99\subset0\ldots249$, byte-identical episodes)---with stable estimates (bootstrap
SE${\approx}0.04$--$0.05$), so the headline---they are calibrated near the oracle and beat the naive
frequentist in the sparse regime but do not exceed a well-posed estimator---is well-supported.  The
open-weight rerun is still accumulating toward $250$, so its rows are noisier; the harness is
model-agnostic and both the skill figure and table accommodate additional cities directly.

\textbf{Information parity.}  We run a parity condition that states the generative structure in the
prompt---the $g$-range, the accumulation, and the noise scales---and read both the behavioural
slope and the model's stated gain (Appendix~\ref{app:exp2-parity-refs}), distinguishing ``cannot
infer'' from ``can use when told,'' against the fair forecaster references of
Appendix~\ref{app:exp2-parity-refs}.  A finer dose-response---vague hint vs.\ a named ``sticky''
pollster vs.\ full disclosure---remains future work.

\textbf{Measurement noise and multi-sampling.}  The four-shock slope is read from single-sample
predictions, so part of the non-monotonicity (Appendix~\ref{app:exp2-robustness}) is sampling noise
that \emph{multi-sample averaging}---querying each shock several times and averaging before taking
the slope---would reduce, tightening the recovery estimates (most in the baseline condition, where
non-monotonicity is highest).  We deliberately do not multi-sample: the monotonic-only control
already shows the plateau is not an attenuation artifact, the primary (parity) condition is already
clean, and a non-trivial share of the inconsistency is \emph{genuine}---a coherent forecaster
should never predict a lower poll for more positive news---so averaging it away would suppress a
behavioral signal that is itself diagnostic of an incoherent world model.  Multi-sampling at the
$\sim$$3\times$ cost remains available for a revision if tighter baseline estimates are required.

\textbf{Calibration vs.\ recovery.}  The over-estimation of $\overline{\hat g}$ by the larger
model suggests a prior that news moves polls more than this world warrants.  Separating the
recovered \emph{structure} from its \emph{magnitude} (e.g.\ allowing a per-model affine
correction) would quantify how much of the residual forecast error is mis-calibration rather than
failure to infer.

\textbf{Continuous vs.\ binary latent.}  We draw $g$ continuously to expose the full recovery
curve; a binary biased/straight design recovers the original two-condition framing as a special
case and could be reported alongside for a cleaner verbal contrast.
